\section*{अध्याय \ref{ch:b1:findim} --- परिमित विमा}

\begin{solution}{exo:b1:findim:1}
\begin{enumerate}
  \item $z = -x - y$: $F = \{(x, y, -x-y)\} =
  \operatorname{Vect}\bigl((1,0,-1), (0,1,-1)\bigr)$; दोनों सदिश \omterm{def:b1:vspaces:free}{स्वतंत्र}
  हैं (\omterm{prop:b1:vspaces:coordinates}{निर्देशांक} देखिए): $\dim F = 2$।
  \item $G = \{(x, x, 2t, t)\} = \operatorname{Vect}\bigl((1,1,0,0),
  (0,0,2,1)\bigr)$: \omterm{def:b1:vspaces:free}{स्वतंत्र}, $\dim G = 2$।
  \item $P(1) = P'(1) = 0$ का अर्थ है $(X-1)^2 \mid P$
  (\cref{prop:b1:poly:multiplicity}): $P = (X-1)^2(aX + b)$। \omterm{def:b1:vspaces:free}{आधार}
  $\bigl((X-1)^2, X(X-1)^2\bigr)$, विमा $2$।
\end{enumerate}
\end{solution}

\begin{solution}{exo:b1:findim:2}
$(3,5,7) = (1,1,1) + 2(1,2,3)$ और $(0,1,2) = (1,2,3) - (1,1,1)$: दोनों पहले
दो सदिशों के संयोजन हैं, और वे \omterm{def:b1:vspaces:free}{स्वतंत्र} हैं (अनुपाती नहीं)। \omterm{def:b1:findim:rank}{कोटि} $2$; \omterm{def:b1:vspaces:span}{जनित
समष्टि} का \omterm{def:b1:vspaces:free}{आधार}: $\bigl((1,1,1), (1,2,3)\bigr)$।
\end{solution}

\begin{solution}{exo:b1:findim:3}
$e_1 = (1,0,0,0)$ और $e_3 = (0,0,1,0)$ जोड़कर देखिए। कुल $\bigl((1,1,0,0),
(0,0,1,1), e_1, e_3\bigr)$ \omterm{def:b1:vspaces:free}{स्वतंत्र} है: शून्य संयोजन $\alpha(1,1,0,0) +
\beta(0,0,1,1) + \gamma e_1 + \delta e_3 = 0$ $(\alpha + \gamma, \alpha,
\beta + \delta, \beta) = 0$ कहता है, अतः $\alpha = \beta = 0$, फिर $\gamma =
\delta = 0$। विमा $4$ में चार \omterm{def:b1:vspaces:free}{स्वतंत्र} सदिश: अतः \omterm{def:b1:vspaces:free}{आधार}
(\cref{prop:b1:findim:twoofthree})।
\end{solution}

\begin{solution}{exo:b1:findim:4}
घातें $0, 1, 2, 3$ युग्मशः भिन्न हैं: अतः \omterm{def:b1:vspaces:free}{स्वतंत्र}
(\cref{prop:b1:vspaces:freecriteria}), और विमा $4$ में चार सदिश: अतः \omterm{def:b1:vspaces:free}{आधार}।
$X^3$ के लिए: नीचे की ओर प्रसार कीजिए,
\[
X(X-1)(X-2) = X^3 - 3X^2 + 2X,
\qquad
X(X-1) = X^2 - X,
\]
अतः $X^3 - X(X-1)(X-2) = 3X^2 - 2X$; और $3X^2 - 2X = 3(X^2 - X) + X =
3\,X(X-1) + X$। इस प्रकार
\[
X^3 = X(X-1)(X-2) + 3\,X(X-1) + 1\cdot X + 0\cdot 1 :
\]
$\bigl(1, X, X(X-1), X(X-1)(X-2)\bigr)$ पर \omterm{prop:b1:vspaces:coordinates}{निर्देशांक} $(0, 1, 3, 1)$। (ये
वेश बदले हुए स्टर्लिंग अंक हैं।)
\end{solution}

\begin{solution}{exo:b1:findim:5}
ग्रासमान: $\dim(F \cap G) = 4 + 5 - \dim(F + G)$, और $F + G$ $E$ की ऐसी
\omterm{def:b1:vspaces:subspace}{उपसमष्टि} है जिसमें $G$ है: $5 \leq \dim(F+G) \leq 7$। अतः $\dim(F \cap G)
\in \{2, 3, 4\}$। $(e_1, \dots, e_7)$ विहित \omterm{def:b1:vspaces:free}{आधार} वाले $\R^7$ में $G =
\operatorname{Vect}(e_1, \dots, e_5)$ लेकर साकार रूप:
\begin{itemize}
  \item $F = \operatorname{Vect}(e_1, e_2, e_6, e_7)$: $F + G = \R^7$,
  सर्वनिष्ठ $\operatorname{Vect}(e_1, e_2)$, विमा $2$;
  \item $F = \operatorname{Vect}(e_1, e_2, e_3, e_6)$: $3$ विमा वाला
  सर्वनिष्ठ;
  \item $F = \operatorname{Vect}(e_1, e_2, e_3, e_4) \subseteq G$: विमा $4$।
\end{itemize}
\end{solution}

\begin{solution}{exo:b1:findim:6}
$H$ एक \omterm{def:b1:vspaces:subspace}{उपसमष्टि} है (\cref{exo:b1:vspaces:1}~(4) अक्षरशः)। गुणनखंड प्रमेय से
(\cref{thm:b1:poly:factor}), $\deg Q \leq n - 1$ के साथ $P \in H \iff P =
(X-1)Q$: \omterm{def:b1:logic:map}{प्रतिचित्रण} $Q \mapsto (X-1)Q$ $\R_{n-1}[X]$ से $H$ पर एक रैखिक
एकैकी आच्छादन है, अतः $H$ का एक \omterm{def:b1:vspaces:free}{आधार} है
\[
\bigl((X-1),\ (X-1)X,\ (X-1)X^2,\ \dots,\ (X-1)X^{n-1}\bigr),
\qquad \dim H = n .
\]
एक \omterm{def:b1:vspaces:sum}{पूरक} रेखा: $\operatorname{Vect}(1)$ (अचर फलन)। वस्तुतः $H \cap
\operatorname{Vect}(1) = \{0\}$ (कोई अशून्य अचर $1$ पर शून्य नहीं होता) और
विमाएँ जुड़कर $n + 1$ देती हैं: अतः ग्रासमान से $H \oplus
\operatorname{Vect}(1) = \R_n[X]$।
\end{solution}

\begin{solution}{exo:b1:findim:7}
विलोपन: $u_p$ हटाने पर \omterm{def:b1:vspaces:span}{जनित समष्टि} केवल सिकुड़ ही सकती है; और वह अधिक से
अधिक एक विमा सिकुड़ती है, क्योंकि छोटी \omterm{def:b1:vspaces:span}{जनित समष्टि} के किसी \omterm{def:b1:vspaces:free}{आधार} में $u_p$
वापस जोड़ने पर बड़ी समष्टि का जनक कुल मिल जाता है, वह भी अधिक से अधिक एक
अतिरिक्त सदिश के साथ। सममित रूप से, कोई सदिश $v$ जोड़ने पर: नई \omterm{def:b1:vspaces:span}{जनित समष्टि}
पुरानी को अधिक से अधिक एक अतिरिक्त जनक के साथ समाहित करती है, अतः उसकी विमा
$r$ है (यदि $v$ पहले से \omterm{def:b1:vspaces:span}{जनित समष्टि} में था) अथवा $r + 1$ (अन्यथा,
\cref{prop:b1:vspaces:freecriteria}~(2) से \omterm{def:b1:vspaces:free}{आधार} बढ़ जाता है)। दोनों \omterm{def:b1:logic:statement}{कथन}
मिलकर ``\omterm{def:b1:findim:rank}{कोटि} $1$-लिप्शिट्ज़ है'' वाला निष्कर्ष देते हैं।
\end{solution}

\begin{solution}{exo:b1:findim:8}
किसी सचमुच वर्धमान शृंखला के अनुदिश विमाएँ सचमुच बढ़ती हैं ($F_i \subseteq
F_{i+1}$, $F_i \neq F_{i+1}$ और \cref{thm:b1:findim:subspaces}: विमाओं की
समानता समष्टियों की समानता बाध्य कर देती)। अतः $\dim F_1 < \dim F_2 < \dots
< \dim F_k$ $\intint{0}{n}$ में पूर्णांकों का सचमुच वर्धमान अनुक्रम है: अधिक
से अधिक $n + 1$ मान, $k \leq n + 1$। $\R^n$ में अधिकतम शृंखला:
\[
\{0\} \subsetneq \operatorname{Vect}(e_1) \subsetneq
\operatorname{Vect}(e_1, e_2) \subsetneq \dots \subsetneq \R^n ,
\]
जिसकी लंबाई ठीक $n + 1$ है।
\end{solution}

\begin{solution}{exo:b1:findim:9}
$F + G$ $F$ को सचमुच समाहित करता है (क्योंकि $G \not\subseteq F$), अतः
$\dim(F + G) = n$ और ग्रासमान $\dim(F \cap G) = (n-1) + (n-1) - n = n - 2$
देता है।

व्यापक दावा, $k$ पर आगमन से: $k$ अधिसमतलों का सर्वनिष्ठ $I_k$ $\dim I_k \geq
n - k$ रखता है। $k = 1$ के लिए सत्य। चरण: $I_{k+1} = I_k \cap H_{k+1}$, और
$E$ के भीतर ग्रासमान:
\[
\dim(I_k \cap H_{k+1}) = \dim I_k + (n - 1) - \dim(I_k + H_{k+1})
\geq \dim I_k + (n-1) - n \geq n - k - 1 . \qedhere
\]
\end{solution}

\begin{solution}{exo:b1:findim:10}
\begin{enumerate}
  \item प्रतिबंध $u_{n+2} - 3u_{n+1} + 2u_n = 0$ रैखिक है और शून्य अनुक्रम
  उसे संतुष्ट करता है: अतः $E$ एक \omterm{def:b1:vspaces:subspace}{उपसमष्टि} है। आगमन से $u_0$ और $u_1$ हर
  $u_n$ तय कर देते हैं, और निर्भरता रैखिक है (हर चरण पिछले दो मानों का रैखिक
  संयोजन है); विलोमतः हर युग्म $(a, b)$ ठीक एक ही $E$ के अनुक्रम से आता है
  ($u_n$ को पुनरावृत्ति से परिभाषित कीजिए)। \cref{ex:b1:findim:dims} की तरह
  $E$ $(u_0, u_1) \in \R^2$ के द्वारा \omterm{def:b1:logic:inj}{एकैकी आच्छादक} तथा रैखिक रूप से
  प्राचलित है: $\dim E = 2$।
  \item अचर: $3\cdot1 - 2\cdot1 = 1$। गुणोत्तर: $3\cdot 2^{n+1} - 2\cdot 2^n
  = (6 - 2)2^n = 2^{n+2}$। दोनों $E$ में हैं। स्वतंत्रता: सभी $n$ के लिए
  $a\cdot 1 + b\cdot 2^n = 0$ से $n = 0$ और $n = 1$ पर मिलता है: $a + b =
  0$, $a + 2b = 0$, अतः $a = b = 0$। विमा $2$ में दो \omterm{def:b1:vspaces:free}{स्वतंत्र} सदिश: अतः \omterm{def:b1:vspaces:free}{आधार}
  (\cref{prop:b1:findim:twoofthree})।
  \item $a + b = 0$, $a + 2b = 1$ हल कीजिए: $b = 1$, $a = -1$, अतः $u_n =
  2^n - 1$ (मर्सेन अनुक्रम)।
\end{enumerate}
\end{solution}

\begin{solution}{exo:b1:findim:11}
\begin{enumerate}
  \item ग्रासमान: $\dim(F \cap G) = \dim F + \dim G - \dim(F + G) \geq \dim
  F + \dim G - n > 0$, अतः $F \cap G \neq \{0\}$। $n = 100$ के साथ: $51 + 50
  - 100 = 1 > 0$, अर्थात् सर्वनिष्ठ में कोई रेखा है।
  \item यदि $F \cap H = \{0\}$, तो योग प्रत्यक्ष है और $\dim F + (n - 1) =
  \dim(F \oplus H) \leq n$, अतः $\dim F \leq 1$। (विलोमतः $H$ में न समाई कोई
  रेखा इसे संतुष्ट करती ही है: अधिसमतल लगभग कुछ भी नहीं चूकते।)
\end{enumerate}
\end{solution}

\begin{solution}{exo:b1:findim:12}
$d = \dim F = \dim G$ पर $d = n$ से $d = 0$ तक नीचे की ओर आगमन। यदि $d = n$:
तो $F = G = E$ और $S = \{0\}$ काम कर जाता है। मान लीजिए \omterm{def:b1:logic:statement}{कथन} $d + 1 \leq n$
विमा वाली उपसमष्टि-जोड़ियों के लिए सत्य है, और $\dim F = \dim G = d < n$
लीजिए।

यदि $F = G$: तो $S$ के लिए $F$ का कोई भी \omterm{def:b1:vspaces:sum}{पूरक} लीजिए
(\cref{thm:b1:findim:subspaces})। यदि $F \neq G$: तो दोनों उचित हैं, अतः
\cref{exo:b1:vspaces:12} से ऐसा $x \notin F \cup G$ है। योग $F' = F \oplus
\operatorname{Vect}(x)$ और $G' = G \oplus \operatorname{Vect}(x)$ प्रत्यक्ष
हैं ($x \notin F$, $x \notin G$) और उनकी विमा $d + 1$ है; आगमन-परिकल्पना से
उनका कोई उभयनिष्ठ \omterm{def:b1:vspaces:sum}{पूरक} $S'$ है: $E = F' \oplus S' = G' \oplus S'$। $S =
\operatorname{Vect}(x) \oplus S'$ रखिए --- यह प्रत्यक्ष है, क्योंकि $S' \cap
\operatorname{Vect}(x) \subseteq S' \cap F' = \{0\}$।

तब $F + S = F + \operatorname{Vect}(x) + S' = F' + S' = E$, और $F \cap S =
\{0\}$: यदि $f \in F$, $s' \in S'$ के साथ $f = \lambda x + s'$, तो $s' = f -
\lambda x \in F' \cap S' = \{0\}$, अतः $f = \lambda x$, जो $\lambda = 0$ ($x
\notin F$) और $f = 0$ बाध्य कर देता है। अतः $E = F \oplus S$, और सममित रूप
से $E = G \oplus S$।
\end{solution}

\begin{solution}{pb:b1:findim:1}
\textbf{1.} $\Q$ में $0 \neq 1$ है, वह योग, गुणन और विपरीत के अंतर्गत स्थायी
है, और हर अशून्य परिमेय का परिमेय प्रतिलोम है: अतः एक \omterm{def:b1:structures:field}{क्षेत्र}
(\cref{def:b1:structures:field})। \omterm{def:b1:vspaces:span}{जनित समष्टि}, स्वतंत्रता और \omterm{def:b1:vspaces:free}{आधार} की
परिभाषाएँ तथा अध्याय 18--19 की उपपत्तियाँ केवल सदिश-योग, वितरण नियम और किसी
\omterm{def:b1:structures:field}{क्षेत्र} के भीतर के अदिश-अंकगणित का प्रयोग करती हैं --- कभी किसी निरपेक्ष
मान, क्रम या सीमा का नहीं। अतः $\R$, अपने स्वयं के योग और गुणन $\Q \times \R
\to \R$ के साथ, एक $\Q$-सदिश समष्टि है, और सभी प्रमेयें लागू होती हैं। अदिश
से भाग \cref{thm:b1:findim:exchange} की उपपत्ति में एक बार आता है: ``$g_k$
के लिए हल करने'' के लिए उसके अशून्य गुणांक से भाग दिया जाता है।

\textbf{2.} यदि $a, b \in \Q$ के दोनों शून्य न हों और $a + b\sqrt2 = 0$: तो
$b \neq 0$ $\sqrt2 = -a/b \in \Q$ देता, जो $\sqrt2$ की अपरिमेयता का विरोध
करता ($p = 2$ के साथ \cref{exo:b1:arith:7}); अतः $b = 0$, फिर $a = 0$:
अर्थात् $\Q$ पर \omterm{def:b1:vspaces:free}{स्वतंत्र}। $\R$ पर संबंध $\sqrt2\cdot 1 + (-1)\cdot\sqrt2 =
0$ अतुच्छ है: अतः परतंत्र। इस प्रकार $\dim_\Q \Q(\sqrt2) = 2$, और \omterm{prop:b1:vspaces:coordinates}{निर्देशांक}
$(a, b)$ अद्वितीय हैं।

\textbf{3.} $(a + b\sqrt2)(c + d\sqrt2) = (ac + 2bd) + (ad + bc)\sqrt2 \in
\Q(\sqrt2)$। तब
\[
\sigma(uv) = (ac + 2bd) - (ad + bc)\sqrt2
= (a - b\sqrt2)(c - d\sqrt2) = \sigma(u)\sigma(v),
\]
अतः $N(uv) = uv\,\sigma(uv) = u\sigma(u)\,v\sigma(v) = N(u)N(v)$। यदि $u
\neq 0$ के साथ $N(u) = a^2 - 2b^2 = 0$: तो $b \neq 0$ $(a/b)^2 = 2$ देता,
अर्थात् $2$ का परिमेय वर्गमूल; अतः $b = 0$, फिर $a = 0$ और $u = 0$:
विरोधाभास। इस प्रकार $u \neq 0$ के लिए $N(u) \neq 0$।

\textbf{4.} $u \cdot \dfrac{\sigma(u)}{N(u)} = \dfrac{N(u)}{N(u)} = 1$, और
$\sigma(u)/N(u) \in \Q(\sqrt2)$: हर अशून्य अवयव $\Q(\sqrt2)$ के भीतर ही
व्युत्क्रमणीय है, अतः वह $\R$ का उपक्षेत्र है। उदाहरण: $N(3 + 2\sqrt2) = 9 -
8 = 1$, अतः
\[
\frac1{3 + 2\sqrt2} = 3 - 2\sqrt2 ;
\qquad
N(1 + \sqrt2) = -1, \quad
\frac1{1 + \sqrt2} = \sqrt2 - 1 .
\]

\textbf{5.} यदि $\sqrt6 = p/q$, तो $6q^2 = p^2$; बाईं ओर $2$ का घातांक $1 +
2v_2(q)$ है, जो विषम है, और दाईं ओर $2v_2(p)$, जो सम है: असंभव। अब मान लीजिए
$a, b \in \Q$ के साथ $\sqrt3 = a + b\sqrt2$। वर्ग करने पर: $3 = a^2 + 2b^2 +
2ab\sqrt2$। यदि $ab \neq 0$, तो $\sqrt2 = (3 - a^2 - 2b^2)/(2ab) \in \Q$:
असंभव। यदि $b = 0$: $\sqrt3 = a \in \Q$, जो \cref{exo:b1:arith:7} ($p = 3$)
का विरोध करता है। यदि $a = 0$: $\sqrt3 = b\sqrt2$, और $\sqrt2$ से गुणा करने
पर: $\sqrt6 = 2b \in \Q$: असंभव। अतः $\sqrt3 \notin \Q(\sqrt2)$।

\textbf{6.} आधार-सदिशों के गुणनफल हैं
\[
\sqrt2\,\sqrt3 = \sqrt6,\quad
\sqrt2\,\sqrt6 = 2\sqrt3,\quad
\sqrt3\,\sqrt6 = 3\sqrt2,\quad
(\sqrt2)^2 = 2,\ (\sqrt3)^2 = 3,\ (\sqrt6)^2 = 6,
\]
और सभी $M$ में हैं। $M$ के दो अवयवों का गुणनफल द्विरैखिकता से इन्हीं के
परिमेय संयोजनों में प्रसारित हो जाता है: अतः $M$ गुणन के अंतर्गत स्थायी है।

\textbf{7.} मान लीजिए $x, y \in K = \Q(\sqrt2)$ के साथ $x + y\sqrt3 = 0$।
यदि $y \neq 0$, तो $\sqrt3 = -x/y \in K$ (प्रश्न 4: $K$ एक \omterm{def:b1:structures:field}{क्षेत्र} है), जो
प्रश्न 5 का विरोध करता है। अतः $y = 0$, फिर $x = 0$: अर्थात् $(1, \sqrt3)$
$K$ पर \omterm{def:b1:vspaces:free}{स्वतंत्र} है। और वह जनक भी है: $K + K\sqrt3 = \{(a + b\sqrt2) + (c +
d\sqrt2)\sqrt3\} = \operatorname{Vect}_\Q (1, \sqrt2, \sqrt3, \sqrt6) = M$।
इस प्रकार $\dim_K M = 2$।

\textbf{8.} संबंध $a + b\sqrt2 + c\sqrt3 + d\sqrt6 = 0$ $K$ में गुणांकों के
साथ $(a + b\sqrt2) + (c + d\sqrt2)\sqrt3 = 0$ के रूप में फिर से समूहबद्ध हो
जाता है; प्रश्न 7 से दोनों शून्य हो जाते हैं, और प्रश्न 2 से $a = b = 0$ तथा
$c = d = 0$। अतः कुल \omterm{def:b1:vspaces:free}{स्वतंत्र} है और $\dim_\Q M = 4$।

\textbf{9.} हर $x \in M$ $y_j \in K$ के साथ $x = \sum_j y_j f_j$ लिखा जाता
है ($K$ पर $M$ का \omterm{def:b1:vspaces:free}{आधार}), और हर $y_j = \sum_i \lambda_{ij} e_i$ $\lambda_{ij}
\in \Q$ के साथ ($\Q$ पर $K$ का \omterm{def:b1:vspaces:free}{आधार}); प्रतिस्थापित करने पर $x = \sum_{i,j}
\lambda_{ij}\, e_i f_j$: अतः गुणनफल $\Q$ पर $M$ को जनित करते हैं।

\textbf{10.} मान लीजिए $\lambda_{ij} \in \Q$ के साथ $\sum_{i,j}
\lambda_{ij}\, e_i f_j = 0$। फिर से समूहबद्ध कीजिए: $\sum_j \bigl(\sum_i
\lambda_{ij} e_i\bigr) f_j = 0$, और भीतरी योग $K$ के सदस्य हैं। $K$ पर
$(f_j)$ की स्वतंत्रता हर $j$ के लिए $\sum_i \lambda_{ij} e_i = 0$ देती है;
फिर $\Q$ पर $(e_i)$ की स्वतंत्रता सभी $i, j$ के लिए $\lambda_{ij} = 0$ दे
देती है।

\textbf{11.} प्रश्न 9 और 10 से $(e_i f_j)_{i \leq m,\, j \leq n}$ $\Q$ पर
$M$ का \omterm{def:b1:vspaces:free}{आधार} है, जिसमें $mn$ अवयव हैं:
\[
\dim_\Q M = m\,n = \dim_\Q K \cdot \dim_K M .
\]
जाँच: $\dim_\Q K = 2$ और $\dim_K M = 2$ (प्रश्न 2 और 7) से $\dim_\Q M = 4$
मिलता है, जो प्रश्न 8 ही है।

\textbf{12.} सदिश $u a_1, \dots, u a_d$ $A$ के सदस्य हैं (स्थायित्व)। वे
\omterm{def:b1:vspaces:free}{स्वतंत्र} हैं: यदि $\sum_i \lambda_i\, u a_i = 0$, तो $\R$ में $u \sum_i
\lambda_i a_i = 0$, और $u \neq 0$ $\sum_i \lambda_i a_i = 0$ बाध्य कर देता
है, अतः $\lambda_i = 0$ ($a_i$ \omterm{def:b1:vspaces:free}{आधार} बनाते हैं)। अतः $(u a_1, \dots, u a_d)$
$A$ में $d$ सदिशों का \omterm{def:b1:vspaces:free}{स्वतंत्र कुल} है, $\dim_\Q A = d$: अतः \omterm{def:b1:vspaces:free}{आधार}
(\cref{prop:b1:findim:twoofthree})। विशेष रूप से $1 \in A$ $v = \sum_i \mu_i
a_i \in A$ के साथ $1 = \sum_i \mu_i\, u a_i = u\,v$ के रूप में विघटित होता
है: अर्थात् $u$ का प्रतिलोम $A$ में है। इससे प्रश्न 4 ($A = \Q(\sqrt2)$) फिर
से मिल जाता है और एक ही झटके में सिद्ध हो जाता है कि $M$ $\R$ का उपक्षेत्र
है (प्रश्न 6 के साथ)।

\textbf{13.} $d + 1$ सदिश $1, x, x^2, \dots, x^{d}$ सभी $A$ में हैं
(गुणनफलों के अंतर्गत स्थायित्व); $A$ के किसी \omterm{def:b1:vspaces:free}{स्वतंत्र कुल} में अधिक से अधिक
$d$ सदिश होते हैं (\cref{prop:b1:findim:twoofthree}), अतः वे परतंत्र हैं:
ऐसे परिमेय $\lambda_0, \dots, \lambda_d$ हैं, जो सब शून्य नहीं हैं, और
$\sum_k \lambda_k x^k = 0$। \omterm{def:b1:poly:def}{बहुपद} $P = \sum_k \lambda_k X^k$ अशून्य है, उसके
गुणांक परिमेय हैं, उसकी घात $\leq d$ है, और $P(x) = 0$।

\textbf{14.} $s^2 = 2 + 2\sqrt6 + 3 = 5 + 2\sqrt6$: \omterm{prop:b1:vspaces:coordinates}{निर्देशांक} $(5, 0, 0,
2)$। तब
\[
s^3 = s\,s^2 = (\sqrt2 + \sqrt3)(5 + 2\sqrt6)
= 5\sqrt2 + 2\sqrt{12} + 5\sqrt3 + 2\sqrt{18}
= 11\sqrt2 + 9\sqrt3 ,
\]
\omterm{prop:b1:vspaces:coordinates}{निर्देशांक} $(0, 11, 9, 0)$ ($\sqrt{12} = 2\sqrt3$, $\sqrt{18} = 3\sqrt2$ का
प्रयोग करते हुए)। अंत में $s^4 = (s^2)^2 = 25 + 20\sqrt6 + 24 = 49 +
20\sqrt6$: \omterm{prop:b1:vspaces:coordinates}{निर्देशांक} $(49, 0, 0, 20)$।

\textbf{15.} मान लीजिए $a + bs + cs^2 + ds^3 = 0$। $(1, \sqrt2, \sqrt3,
\sqrt6)$ पर चारों \omterm{prop:b1:vspaces:coordinates}{निर्देशांक} पढ़ने पर:
\[
a + 5c = 0,\qquad b + 11d = 0,\qquad b + 9d = 0,\qquad 2c = 0 .
\]
अतः $c = 0$, फिर $a = 0$; बीच के समीकरणों को घटाने पर $2d = 0$, फिर $b = 0$:
अतः $(1, s, s^2, s^3)$ \omterm{def:b1:vspaces:free}{स्वतंत्र} है। $4$ विमा वाले $M$ में चार \omterm{def:b1:vspaces:free}{स्वतंत्र} सदिश:
अतः \omterm{def:b1:vspaces:free}{आधार}, इसलिए $\operatorname{Vect}_\Q(1, s, s^2, s^3) = M$। प्रश्न 14 से
$s^3 - 9s = 2\sqrt2$ और $11s - s^3 = 2\sqrt3$:
\[
\sqrt2 = \frac{s^3 - 9s}{2},
\qquad
\sqrt3 = \frac{11s - s^3}{2} .
\]

\textbf{16.} $s^4 - 10s^2 + 1 = (49 + 20\sqrt6) - 10(5 + 2\sqrt6) + 1 = 0$।
$s$ पर शून्य होने वाला $\leq 3$ घात का कोई \omterm{def:b1:poly:def}{इकाई-अग्र बहुपद} $(1, s, s^2,
s^3)$ का अतुच्छ शून्य संयोजन दे देता, जो प्रश्न 15 का विरोध करता: अतः $X^4 -
10X^2 + 1$ की घात न्यूनतम है। उसके मूल: $X^2 = 5 \pm 2\sqrt6 = (\sqrt3 \pm
\sqrt2)^2$, अतः चारों वास्तविक मूल $\pm(\sqrt3 + \sqrt2)$ और $\pm(\sqrt3 -
\sqrt2)$ हैं, अर्थात् $\pm\sqrt2 \pm \sqrt3$।

\textbf{17.} $s^4 - 10s^2 + 1 = 0$ से: $s\,(10s - s^3) = 1$, अतः
\[
\frac1s = 10s - s^3 = 10(\sqrt2 + \sqrt3) - (11\sqrt2 + 9\sqrt3)
= \sqrt3 - \sqrt2 ,
\]
जो $(\sqrt3 + \sqrt2)(\sqrt3 - \sqrt2) = 1$ के संगत है। यदि $a + b\sqrt2 +
c\sqrt3 + d\sqrt6 = r \in \Q$, तो $(a - r) + b\sqrt2 + c\sqrt3 + d\sqrt6 =
0$, और स्वतंत्रता (प्रश्न 8) $b = c = d = 0$ बाध्य कर देती है (तथा $a = r$)।
$\sqrt2 + \sqrt3 + \sqrt6$ के लिए \omterm{prop:b1:vspaces:coordinates}{निर्देशांक} $(0, 1, 1, 1)$ इस रूप के नहीं
हैं: अतः अपरिमेय --- यह \omterm{def:b1:logic:statement}{कथन} \cref{exo:b1:reals:7} से सुदृढ़ है, और वह भी
वर्ग करने की किसी चाल के बिना प्राप्त।

\textbf{18.} $X^3 - 2$ का कोई परिमेय मूल $p/q$ (लघुतम पदों में) $p^3 = 2q^3$
संतुष्ट करता है, और \cref{exo:b1:poly:5} की कसौटी $p \mid 2$, $q \mid 1$
देती है: अतः संभावित मान $\pm1, \pm2$, जिनके घन $\pm1, \pm8 \neq 2$ हैं। कोई
परिमेय मूल नहीं; विशेष रूप से $t = 2^{1/3} \notin \Q$, अतः $(1, t)$ $\Q$ पर
\omterm{def:b1:vspaces:free}{स्वतंत्र} है (प्रश्न 2 की तरह)।

\textbf{19.} मान लीजिए $(1, t, t^2)$ परतंत्र है: $\deg P \leq 2$ के साथ कोई
अशून्य $P \in \Q[X]$ ऐसा है कि $P(t) = 0$; ऐसा $P$ \emph{न्यूनतम} घात $d
\geq 1$ वाला चुनिए। प्रश्न 18 से $d \neq 1$, अतः $d = 2$। यूक्लिडीय भाग
(\cref{thm:b1:poly:division}): $Q, R \in \Q[X]$, $\deg R < 2$ के साथ $X^3 -
2 = PQ + R$। $t$ पर मूल्यांकन करने पर: $0 = P(t)Q(t) + R(t) = R(t)$, अतः $R$
$t$ पर शून्य होता है; $d$ की न्यूनतमता $R = 0$ बाध्य कर देती है। तब $\deg Q
= 1$ के साथ $X^3 - 2 = PQ$: $Q$ का परिमेय मूल $X^3 - 2$ का परिमेय मूल है, जो
प्रश्न 18 का विरोध करता है। अतः $(1, t, t^2)$ \omterm{def:b1:vspaces:free}{स्वतंत्र} है और $\dim_\Q A =
3$। स्थायित्व: $t^3 = 2$ $1, t, t^2$ के हर गुणनफल को उन्हीं के संयोजन तक घटा
देता है ($t\cdot t^2 = 2$, $t^2\cdot t^2 = 2t$); $A$ में $1$ है: अतः प्रश्न
12 से $A$ $\R$ का उपक्षेत्र है।

\textbf{20.} मान लीजिए $t \in M$। तब $t^2 \in M$ (स्थायित्व), अतः $A
\subseteq M$, और $M$ \omterm{def:b1:structures:field}{क्षेत्र} $A$ पर एक \omterm{def:b1:vspaces:def}{सदिश समष्टि} है: अभिगृहीत वही
$\R$-अंकगणित के हैं, बस सीमित करके। वह $A$ पर \omterm{def:b1:findim:def}{परिमित-विमीय} है (कोई परिमित
$\Q$-जनक कुल $A \supseteq \Q$ पर तो और भी सुगमता से जनक है)। $\Q \subseteq A
\subseteq M$ के लिए \omterm{pb:b1:findim:1}{मीनार नियम} देता है
\[
4 = \dim_\Q M = \dim_\Q A \cdot \dim_A M = 3\,\dim_A M ,
\]
जो असंभव है: $3$ $4$ को विभाजित नहीं करता। अतः $2^{1/3} \notin M$: $1,
\sqrt2, \sqrt3, \sqrt6$ का कोई परिमेय संयोजन $2^{1/3}$ के बराबर नहीं है।

\textbf{21.} $\Q(\sqrt2) \subseteq M$, अतः $t \notin \Q(\sqrt2)$। यदि $t +
\sqrt2 = r \in \Q$, तो $t = r - \sqrt2 \in \Q(\sqrt2)$: विरोधाभास ---
अर्थात् $2^{1/3} + \sqrt2$ अपरिमेय है। यदि $t = a + b\sqrt3$, तो $t \in
\operatorname{Vect}_\Q(1, \sqrt3) \subseteq M$: फिर से विरोधाभास।

\textbf{22.} $B$ \omterm{def:b1:structures:field}{क्षेत्र} $A$ पर एक \omterm{def:b1:vspaces:def}{सदिश समष्टि} है (अदिशों का सीमन), और
\omterm{def:b1:findim:def}{परिमित-विमीय} है, क्योंकि $\dim_\Q B$ परिमित है। $\Q \subseteq A \subseteq B$
के लिए \omterm{pb:b1:findim:1}{मीनार नियम} $\dim_\Q B = \dim_\Q A \cdot \dim_A B$ देता है: अर्थात्
बायाँ गुणनखंड \omterm{def:b1:arith:divides}{विभाजित करता है}। अतः $M$ के उपक्षेत्रों की $\Q$-विमा $1$, $2$
या $4$ है: विमा $1$ स्वयं $\Q$ है, और विमा $4$ $M$ है।

\textbf{23.} मान लीजिए $F \subseteq M$ ऐसा उपक्षेत्र है कि $\dim_\Q F = 2$,
और $w \in F \setminus \Q$: $(1, w)$ \omterm{def:b1:vspaces:free}{स्वतंत्र} है, अतः $F$ का \omterm{def:b1:vspaces:free}{आधार}। प्रश्न 13
से ($d = 2$) परिमेय $\alpha, \beta$ के लिए $w^2 = \alpha + \beta w$। $v = w
- \beta/2 \in F \setminus \Q$ रखने पर:
\[
v^2 = w^2 - \beta w + \frac{\beta^2}4 = \alpha + \frac{\beta^2}4
\;\in\; \Q,
\]
और $F = \operatorname{Vect}(1, v)$। अब $v = a + b\sqrt2 + c\sqrt3 + d\sqrt6$
लिखिए और प्रसार कीजिए:
\[
v^2 = (a^2 + 2b^2 + 3c^2 + 6d^2)
+ (2ab + 6cd)\sqrt2 + (2ac + 4bd)\sqrt3 + (2ad + 2bc)\sqrt6 .
\]
तीनों अपरिमेय \omterm{prop:b1:vspaces:coordinates}{निर्देशांक} शून्य हो जाते हैं: $ab + 3cd = ac + 2bd = ad + bc =
0$। यदि $a \neq 0$: तो $b = -3cd/a$, और दूसरी शर्त $c(a^2 - 6d^2) = 0$ बन
जाती है; चूँकि $d \neq 0$ के साथ $a^2 = 6d^2$ $\sqrt6 = \abs{a/d}$ को परिमेय
बना देता, अतः या तो $c = 0$ या $d = 0$, और दोनों स्थितियों में शेष शर्तें $b
= c = d = 0$ बाध्य कर देती हैं, अर्थात् $v \in \Q$: जो निषिद्ध है। अतः $a =
0$, और शर्तें $3cd = 2bd = bc = 0$ कहती हैं: $b, c, d$ में से अधिक से अधिक
एक अशून्य है। इस प्रकार $v$ $\sqrt2$, $\sqrt3$ अथवा $\sqrt6$ का परिमेय गुणज
है, और
\[
F \in \bigl\{\Q(\sqrt2),\ \Q(\sqrt3),\ \Q(\sqrt6)\bigr\} ,
\]
और इनमें से हर एक सचमुच $2$-विमीय उपक्षेत्र है (स्थायित्व प्रश्न 6 की तरह,
प्रतिलोम प्रश्न 12 से): अतः ठीक तीन द्विघातीय उपक्षेत्र।

\textbf{24.} $(1 + \sqrt2 + \sqrt3)(1 + \sqrt2 - \sqrt3) = (1 + \sqrt2)^2 -
3 = 2\sqrt2$, अतः
\[
\frac1{1 + \sqrt2 + \sqrt3}
= \frac{1 + \sqrt2 - \sqrt3}{2\sqrt2}
= \frac{\sqrt2 + 2 - \sqrt6}{4}
= \frac12 + \frac14\sqrt2 + 0\cdot\sqrt3 - \frac14\sqrt6 .
\]
(जाँच: प्रसार के बाद $(1 + \sqrt2 + \sqrt3)(2 + \sqrt2 - \sqrt6) = 4$।)

\textbf{25.} (क) $\Q$-विमा $\R$ के हर उपक्षेत्र से एक \emph{पूर्णांक} जोड़
देती है, और \omterm{pb:b1:findim:1}{मीनार नियम} इन पूर्णांकों को अंतर्विष्टताओं के अनुदिश गुणा करा
देता है: विमा किसी अंकगणितीय अचर की तरह व्यवहार करती है, और \omterm{def:b1:arith:divides}{विभाज्यता} की
बाध्यताएँ असंभवता की उपपत्तियाँ बन जाती हैं। (ख) \omterm{def:b1:findim:def}{परिमित विमा} हर अवयव को ऐसा
अशून्य परिमेय बहुपद-समीकरण संतुष्ट करने पर बाध्य कर देती है जिसकी घात अधिक
से अधिक विमा जितनी हो: परिमितता का अर्थ बीजीयता है। (ग) तीन में से दो वाले
नियम ने ``$u$ से गुणा किसी \omterm{def:b1:vspaces:free}{आधार} को अधिकतम आकार के \omterm{def:b1:vspaces:free}{स्वतंत्र कुल} पर भेजता है''
को आच्छादकता में बदल दिया, और बिना किसी सूत्र के $u^{-1}$ पैदा कर दिया:
प्रतिलोम गिनती से आए। (घ) \omterm{pb:b1:findim:1}{मीनार नियम} $4$-विमीय \omterm{def:b1:structures:field}{क्षेत्र} के भीतर $3$-विमीय
\omterm{def:b1:structures:field}{क्षेत्र} को मना करता है, और इसीलिए $2^{1/3}$ तक $\sqrt2$ और $\sqrt3$ से नहीं
पहुँचा जा सकता। भाग III की प्रमेय \emph{डेडेकिंड का \omterm{pb:b1:findim:1}{मीनार नियम}} है, जो गैलुआ
सिद्धांत की पहली चाल है और स्नातक वर्ष~3 के खंड में विकसित होती है।
\end{solution}
